% =====================================================================
%  TITRE II : DES PERSONNELS DU DÉPARTEMENT
% =====================================================================
\titrepartie{Titre II}{Des personnels du département}

% ---------------------------------------------------------------------
\article{Tenue, apparence et attitude}

\chapeau{Le présent article s'applique à l'ensemble des personnels. Il constitue une condition d'emploi, tant pour les candidats que pour les agents en fonction. La non-conformité à ses dispositions peut constituer un motif de révocation.}

\cl[apparence]{Tout agent conserve en tout temps une apparence et une attitude professionnelles, ainsi qu'une hygiène convenable. L'apparence de l'agent est le premier élément de crédibilité du département auprès du public.}

\sstitre{Coiffure}

\cl{\stress{Hommes} : le sommet de la tête est soigneusement coiffé. Les cheveux présentent un aspect effilé et, une fois peignés, ne dépassent ni les oreilles ni les sourcils. La barbe est rasée ou taillée court et régulièrement entretenue.}

\cl{\stress{Femmes} : les cheveux sont coiffés de manière à ne gêner ni la vision ni les oreilles, et à permettre le port de tout couvre-chef du service. Ils sont de préférence attachés pendant toute la durée du service. Le maquillage est autorisé dans des teintes naturelles, à l'exclusion des maquillages pailletés, sombres ou fluorescents.}

\cl{Les colorations capillaires, quelles qu'elles soient, sont interdites.}

\sstitre{Tatouages et bijoux}

\cl{Les tatouages sont interdits sur le visage et sur les mains. Ils sont autorisés sur les bras, le buste et les jambes, et peuvent être apparents, sous réserve de ne représenter aucune image offensive, discriminatoire, morbide ou associée à une organisation criminelle.}

\cl{Les boucles d'oreilles et les colliers apparents sont interdits. Les autres bijoux (bracelet, montre, alliance) sont autorisés dès lors qu'ils demeurent discrets et ne présentent aucun risque en intervention.}

\sstitre{Hygiène opérationnelle}

\cl[gants]{L'utilisation de gants en latex bleu est obligatoire pour toute palpation, fouille ou manipulation d'effets personnels. L'agent en dispose en permanence sur lui.}

\begin{note}[Adaptation en service d'investigation]
Les agents affectés à une mission d'infiltration ou de surveillance discrète peuvent être exemptés des règles d'apparence du présent article pour la durée de la mission, sur note de service du Command Staff précisant la nature et la durée de la dérogation.
\end{note}

% ---------------------------------------------------------------------
\article{Uniformes et insignes}

\cl[uniformes]{Les personnels portent exclusivement les uniformes fournis par le département. Les tenues réservées aux divisions portent la mention de la division concernée ; leur port est strictement réservé aux membres confirmés de celle-ci.}

\cl{Les Police Officer I probatoires portent, jusqu'à la validation de leur First Lincoln, une tenue à manches longues avec cravate. Le port du képi peut leur être imposé.}

\cl[tenue-civile]{Les Detectives en enquête et les membres du Command Staff sont exemptés de l'obligation de port de l'uniforme et peuvent revêtir une tenue civile. Lorsqu'ils agissent en qualité d'agent du département, ils portent de façon apparente et lisible soit la plaque du service, soit un gilet pare-balles portant la mention « Police ».}

\cl{D'autres agents peuvent être exemptés à la discrétion du Command Staff. Cette exemption fait l'objet d'une note de service précisant les éléments qui la justifient et sa durée.}

\cl{Tous les trente jours de service accomplis, l'agent est autorisé à porter une barrette d'ancienneté sur sa manche, dans la limite de six barrettes. Le décompte figure sur le terminal de l'agent.}

\cl{La plaque et le numéro de badge sont portés de manière apparente et lisible en toute circonstance de service en uniforme. Dissimuler, masquer ou retirer son identification en intervention constitue une faute grave.}

\sstitre{Insignes de grade}

\vspace{2pt}
\noindent
\begin{tabular}{C{20mm} P{40mm} P{84mm}}
\headrow
\thead{Insigne} & \thead{Grade} & \thead{Port} \\
\insigne{images/g-chief.png} & \tkey{Chief} & \tcell{Quatre étoiles, col de chemise.} \\
\altrow
\insigne{images/g-commander.png} & \tkey{Commander} & \tcell{Étoile, col de chemise.} \\
\insigne{images/g-captain.png} & \tkey{Captain} & \tcell{Double barrette, col de chemise.} \\
\altrow
\insigne{images/g-lieutenant.png} & \tkey{Lieutenant} & \tcell{Barrette simple, col de chemise.} \\
\insigne{images/g-sergeant2.png} & \tkey{Sergeant II} & \tcell{Trois chevrons et un arceau, manche.} \\
\altrow
\insigne{images/g-sergeant1.png} & \tkey{Sergeant I} & \tcell{Trois chevrons, manche.} \\
\insigne{images/g-slo.png} & \tkey{Senior Lead Officer} & \tcell{Deux chevrons et une étoile, manche.} \\
\altrow
\insigne{images/g-officer3.png} & \tkey{Police Officer III} & \tcell{Deux chevrons, manche.} \\
\bottomrule
\end{tabular}

\vspace{3pt}
{\lspdsans\fontsize{8}{10}\selectfont\color{lspdgrey}
Les grades de Police Officer I et II ne comportent pas d'insigne de manche. Les grades de Detective portent l'insigne du grade de terrain équivalent.}

% ---------------------------------------------------------------------
\article{Équipement et armement}

\cl[equipement]{Les personnels n'utilisent que l'équipement fourni par le département. En service, chaque agent s'assure de disposer en permanence de son équipement complet et opérationnel. À chaque fin de service, l'intégralité de l'équipement est déposée dans les locaux du département. Tout manquement à cette règle expose l'agent à une révocation.}

\cl{En cas de perte, de vol ou de destruction d'un élément d'équipement, l'agent le signale sans délai à sa hiérarchie et rédige un rapport circonstancié. L'absence de signalement est sanctionnée indépendamment de la perte elle-même.}

\cl{Les superviseurs et la direction de la Police Academy contrôlent l'état des stocks et informent le Command Staff des besoins en équipement.}

\sstitre{Dotation individuelle}

\begin{multicols}{2}
\begin{itemize}[leftmargin=5mm,topsep=0pt]
  \item 1 pistolet Glock 17
  \item 1 taser Axon T7
  \item 2 paires de menottes
  \item 1 clé de menottes
  \item 1 matraque télescopique
  \item 1 lampe torche
  \item 1 radio de service
  \item 1 gilet pare-balles
  \item 1 caméra-piéton, intégrée à la tenue
  \item Gants d'intervention en latex
\end{itemize}
\end{multicols}

\cl{Les agents titulaires des accréditations correspondantes peuvent se voir remettre des matériels complémentaires.}

\sstitre{Port de l'arme hors service}

\cl[arme-hors-service]{Tout agent titulaire au minimum du grade de Police Officer III est autorisé à conserver hors service une arme de poing, un chargeur, des liens de contrainte et une radio. L'arme demeure impérativement dissimulée et n'est jamais exhibée en public ; le port d'un holster apparent est interdit. Son emploi est strictement réservé à la légitime défense.}

\cl{Tout emploi de l'arme hors service donne lieu, sans exception, à un rapport détaillé exposant les circonstances de l'incident et les motifs du recours à l'arme.}

\cl{En cas d'abus ou d'usage inapproprié, le Command Staff ou l'État-Major peut retirer à l'agent l'autorisation de port hors service, par décision motivée.}

\sstitre{Habilitations d'armement}

\vspace{2pt}
\noindent
\begin{tabular}{P{44mm} P{38mm} P{62mm}}
\headrow
\thead{Matériel} & \thead{Grade minimal} & \thead{Condition} \\
\tcell{Lanceur à balles de défense} & \tcell{Police Officer III} & \tcell{Formation interne suivie et validée.} \\
\altrow
\tcell{Fusil à pompe (PPA II)\newline Pistolet-mitrailleur (PPA III)} & \tcell{Senior Lead Officer} & \tcell{Formation interne suivie et validée.} \\
\tcell{Fusil d'assaut (PPA IV)} & \tcell{Sergeant I} & \tcell{Formation interne suivie et validée ; habilitations précédentes acquises.} \\
\bottomrule
\end{tabular}

\cl{Le Command Staff et l'État-Major peuvent retirer à un agent une habilitation d'armement lorsque celle-ci a été employée de manière inappropriée ou contraire à l'intérêt des missions assignées. Le retrait est motivé et notifié par écrit.}

\cl[lbd]{L'emploi du lanceur à balles de défense est strictement réservé aux situations d'émeute et aux cas où un groupe d'individus se dirige de manière agressive ou menaçante vers un agent. Tout emploi en dehors de ces circonstances constitue une violation des protocoles.}

\begin{interdit}[Équipement non réglementaire]
Il est interdit d'emporter en service une arme, une munition, un moyen de contrainte ou un dispositif quelconque non fourni par le département, y compris s'il est détenu légalement à titre personnel. La découverte d'un équipement non réglementaire sur un agent en service entraîne son retrait immédiat et une procédure disciplinaire.
\end{interdit}

% ---------------------------------------------------------------------
\article{Conduite hors service, image et confidentialité}

\cl[hors-service]{Même hors service, l'agent représente le département. Il lui est interdit de s'engager dans tout acte délictuel ou criminel. Le cas échéant, il s'expose à une sanction disciplinaire indépendamment des suites judiciaires.}

\cl{L'agent condamné pour un délit ou un crime est révoqué. Le casier vierge est une condition d'emploi permanente.}

\cl{Hors service, l'agent ne dispose d'aucune prérogative de police. Confronté à d'autres agents en intervention, il s'identifie immédiatement comme agent hors service et se conforme à leurs instructions.}

\cl{Il est interdit de participer, de quelque manière que ce soit, à une opération de police hors service. Cette interdiction s'applique sans exception aux agents suspendus. La légitime défense et l'assistance à personne en péril immédiat demeurent réservées.}

\cl{Les véhicules de service ne peuvent être utilisés hors service, à l'exception des véhicules banalisés attribués nominativement aux enquêteurs et aux divisions autorisées, dans les conditions fixées par note de service.}

\sstitre{Image du département et communication}

\cl[reseaux]{Il est interdit de publier ou de diffuser, sur quelque support que ce soit, des photographies de scène, des images de détenus ou de victimes, des pièces de dossier ou des informations relatives à une enquête en cours.}

\cl{Il est interdit de tenir publiquement des propos discriminatoires, violents ou dénigrant le département, ses membres ou ses décisions, y compris depuis un compte personnel identifiable. La critique interne se formule par les voies prévues à cet effet.}

\cl[presse]{Seuls le Chief ou le porte-parole qu'il désigne s'expriment au nom du département auprès de la presse. Tout agent sollicité par un journaliste renvoie la demande vers le Command Staff.}

\sstitre{Accès aux fichiers}

\cl[fichiers]{La consultation des fichiers du département (terminal embarqué, casiers, adresses, immatriculations, mains courantes) n'est autorisée que pour un motif professionnel identifiable et documenté.}

\begin{interdit}[Consultation personnelle]
Il est interdit de consulter un fichier du département au sujet de soi-même, d'un proche, d'un collègue ou de toute personne dans un cadre autre que professionnel. Cette faute est lourde : elle entraîne la révocation et la saisine de l'autorité judiciaire.
\end{interdit}

% ---------------------------------------------------------------------
\article{Rémunération, frais et primes}

\cl[remuneration]{L'ensemble des personnels perçoit une rémunération hebdomadaire, versée en principe chaque dimanche. Les salaires sont calculés en fonction du grade occupé et intègrent, le cas échéant, les notes de frais et indemnités validées au titre de la période concernée.}

\cl{Le financement du département est assuré par des subventions gouvernementales garantissant la stabilité économique du service et la régularité des versements. Des retards exceptionnels peuvent survenir en cas de contrainte administrative ou de délai bancaire indépendants du service de comptabilité ; les agents concernés en sont informés dans les meilleurs délais.}

\sstitre{Grille salariale hebdomadaire}

\vspace{2pt}
\noindent
\begin{tabular}{P{62mm} P{30mm} P{50mm}}
\headrow
\thead{Grade} & \thead{Montant} & \thead{Observations} \\
\tkey{Police Officer I probatoire} & \tcell{1\,000 \$} & \tcell{Période d'essai.} \\
\altrow
\tkey{Police Officer I confirmé}   & \tcell{1\,000 \$} & \tcell{} \\
\tkey{Police Officer II}   & \tcell{1\,100 \$} & \tcell{} \\
\altrow
\tkey{Police Officer III}  & \tcell{1\,200 \$} & \tcell{} \\
\tkey{Senior Lead Officer} & \tcell{1\,300 \$} & \tcell{Detective I.} \\
\altrow
\tkey{Sergeant I}          & \tcell{1\,400 \$} & \tcell{Detective II.} \\
\tkey{Sergeant II}         & \tcell{1\,500 \$} & \tcell{Detective III.} \\
\altrow
\tkey{Lieutenant}          & \tcell{1\,600 \$} & \tcell{} \\
\tkey{Captain}             & \tcell{1\,700 \$} & \tcell{} \\
\altrow
\tkey{Commander}           & \tcell{1\,800 \$} & \tcell{} \\
\tkey{Chief}               & \tcell{1\,900 \$} & \tcell{} \\
\altrow
\tkey{Secrétaire}          & \tcell{1\,300 \$} & \tcell{Personnel administratif.} \\
\bottomrule
\end{tabular}

\vspace{3pt}
{\lspdsans\fontsize{8}{10}\selectfont\color{lspdgrey}
Les grades de Detective sont rémunérés au niveau du grade de terrain équivalent.}

\cl{Tout versement de prime exceptionnelle est validé par le Command Staff. Toute révision de la grille salariale relève de l'État-Major, conformément aux politiques budgétaires en vigueur.}

\sstitre{Notes de frais et primes}

\cl[frais]{Les agents peuvent bénéficier du remboursement des dépenses engagées dans l'exercice direct de leurs fonctions. Ces dépenses sont justifiées par une facture ou un reçu officiel et validées par la hiérarchie avant leur intégration à la paie hebdomadaire.}

\cl{L'État-Major, le Command Staff et le service de comptabilité peuvent refuser un remboursement. Ce refus est motivé de manière claire et argumentée, et peut faire l'objet d'un recours.}

\cl{Des primes exceptionnelles peuvent être attribuées en reconnaissance d'actes remarquables ou de missions à risque. Elles sont octroyées et validées par le Command Staff, dans la limite du budget alloué.}

\cl[retenue]{Tout comportement contraire au présent règlement, toute faute professionnelle et tout manquement grave à l'éthique peuvent entraîner la suspension partielle ou totale de la rémunération variable, ainsi qu'une retenue sur salaire en cas de dégradation volontaire de matériel appartenant au département.}

\cl{Les agents font preuve de rigueur dans la gestion de leurs demandes financières, dans un souci de transparence et de respect des fonds publics alloués au service.}
